\noindent\textbf{(a)} La respuesta parte de $(0,0)$ con tangente horizontal, supera el equilibrio final $3/16\,\mathrm m$ y oscila acercándose a él. Su primer máximo ocurre en $t=5\pi/\sqrt{19}\simeq3.60365\,\mathrm s$ y vale $\frac3{16}(1+e^{-\pi/\sqrt{19}})\simeq0.278699\,\mathrm m$.

\noindent\textbf{(b)} Es una respuesta subamortiguada alrededor del nuevo equilibrio $15/80=3/16\,\mathrm m$. Sus máximos decrecen y sus mínimos crecen hacia ese valor; no es una aproximación monótona.
